% Part: methods
% Chapter: proofs
% Section: reading-proofs

\documentclass[../../../include/open-logic-section]{subfiles}

\begin{document}

\olfileid{mth}{prf}{rea}
\olsection{د ثبوتونو لوستل}

په درسي کتابونو او مقالو کښې موندل کېدونکي ثبوتونه ډېر کم هغه
ټول تفصيلات ليکي چې موږ تر اوسه په خپلو بېلګو کښې شامل کړي دي۔
ليکوال ډېر ځله نۀ څرګندوي چې کله حالتونه بېلوي، کله نامستقيم
ثبوت ورکوي، يا دا هم نۀ يادوي چې يو تعريف کاروي۔ نو په درسي
کتاب کښې د ثبوت د لوستلو پر وخت به ډېر ځله د هغه د پوهېدو
دپاره دا تفصيلات پخپله بشپړوئ۔ دا کار د هغو بېلابېلو ګامونو
د زده کولو ښه تمرين هم دے چې په ثبوت کښې ئې اخيستل پکار دي۔
يوه بېلګه ګورو۔

\begin{prop}[جذب]
د ټولو سټونو $A$، $B$ دپاره:
\[
A \cap (A \cup B) = A
\]
\end{prop}

\begin{proof}
که $z \in A \cap (A \cup B)$ وي، نو $z \in A$؛ له دې امله
$A \cap (A \cup B)
\subseteq A$۔ اوس $z \in A$ فرض کړئ۔ بيا $z \in A \cup B$ هم
دے، او ځکه $z \in A \cap (A \cup B)$ هم دے۔
\end{proof}

د جذب د قانون پورتنے ثبوت ډېر لنډ شوے دے۔ نۀ د کارول شوو
تعريفونو يادونه لري، نۀ د يوې خبرې تر ثابتولو وړاندې «دا
ثابتول پکار دي» وايي، او داسې نور۔ راځئ پرانيزو ئې۔ ثابته
شوې قضيه د هر ډول سټونو $A$ او $B$ په اړه عمومي دعوه ده؛
نو کله چې ثبوت $A$ يا $B$ يادوي، دا د اختياري سټونو متغيرونه
دي۔ ثبوت هغه عمومي دعوې ثابتوي چې د سټونو د برابرۍ دپاره
اړينې دي: يعنې د برابرۍ د چپ اړخ هر !!{element} د ښي اړخ
!!a{element} دے، او برعکس۔

\begin{quote}
«که $z \in A \cap (A \cup B)$ وي، نو $z \in A$؛ له دې امله
  $A \cap (A \cup B)
  \subseteq A$۔»
\end{quote}

دا د برابرۍ د ثبوت لومړے نيم دے: ثابتوي چې که اختياري~$z$
د چپ اړخ !!a{element} وي، نو د ښي اړخ !!a{element} هم دے؛
يعنې $A \cap (A \cup B) \subseteq A$۔ فرض کړئ
$z \in A \cap (A \cup B)$۔ څرنګه چې $z$ د دوو سټونو د تقاطع
!!{element} هغه وخت او يوازې هغه وخت دے چې د دواړو سټونو
!!{element} وي، نو $z \in A$ او $z \in A \cup B$ دواړه اخلو۔
په ځانګړي ډول $z \in A$، چې همدا مو ثابتول غوښتل۔ څرنګه چې
د لومړي نيم دپاره همدومره کار پکار دے، پوهېږو چې پاتې ثبوت
بايد د دويم نيم ثبوت وي؛ يعنې دا ثابت کړي چې
$A \subseteq A \cap
(A \cup B)$۔

\begin{quote}
«اوس $z \in A$ فرض کړئ۔ بيا $z \in A \cup B$ هم دے،
او ځکه $z \in A \cap (A \cup B)$ هم دے۔»
\end{quote}

په $z \in A$ فرضولو پيل کوو، ځکه ښيو چې د هر~$z$ دپاره
که $z \in A$ وي، نو $z \in A \cap (A \cup B)$۔ د دې
د ښودلو دپاره چې $z
\in A \cap (A \cup B)$، د «$\cap$» د تعريف له مخې بايد
دواړه وښيو: (i) $z \in A$ او (ii) $z \in A \cup B$۔
دلته (i) هماغه زموږ فرض دے؛ نو نور ثبوت ته اړتيا نۀ لري،
او همدا وجه ده چې ثبوت ئې بيا نۀ يادوي۔ د (ii) دپاره
راياد کړئ چې $z$ د سټونو د اتحاد !!a{element} هغه وخت او
يوازې هغه وخت دے چې لږ تر لږه د يوه سټ !!{element} وي۔
څرنګه چې $z \in A$ او $A \cup B$ د $A$ او $B$ اتحاد دے،
دا شرط دلته پوره دے۔ نو $z \in A \cup B$۔ موږ دواړه
ښودلي دي: (i) $z \in A$ او (ii) $z \in A \cup B$؛ ځکه
د «$\cap$» له تعريفه $z \in A \cap (A \cup B)$۔
ثبوت دغه تعريفونه نۀ يادوي؛ فرض شوې ده چې لوستونکے له
مخکې پرې ښه پوهېږي۔ که تاسو پرې نۀ پوهېږئ، بايد بېرته
ورته لاړ شئ او راياد ئې کړئ۔ بيا به دا هم پېژنئ چې ولې
له $z
\in A$ څخه $z \in A \cup B$ راځي، او له $z \in A$ او
$z \in A \cup B$ څخه $z \in A \cap (A \cup B)$ راځي۔

دا د پورتني ثبوت بله بڼه ده، چې هر څه پکښې ښکاره شوي دي:
\begin{proof}{}
[د سټونو د $=$ له تعريفه، د $A \cap (A \cup B) = A$ دپاره
بايد (a) $A \cap (A \cup B) \subseteq A$ او
(b) $A \cap (A \cup B)
  \subseteq A$ وښيو۔ \emph{د سرچينې يادونه (OLSTH-040):
د «ب» په چاپي عبارت کښې لومړے شمول بيا ليکل شوے دے؛
دويم نيم د مقابل لوري شمول غواړي، لکه لاندې تشريح چې ئې
ښيي۔ اصلي فورمول ساتل شوے دے۔} (a): د $\subseteq$ له تعريفه
بايد وښيو چې که $z \in A \cap (A \cup B)$ وي، نو
$z \in A$۔] که $z \in A
\cap (A \cup B)$ وي، نو $z \in A$ [ځکه د $\cap$ له تعريفه
$z
  \in A \cap (A \cup B)$ هغه وخت او يوازې هغه وخت چې
$z \in A$ او $z \in A \cup B$]؛ نو
$A
\cap (A \cup B) \subseteq A$۔ [(b): د $\subseteq$ له تعريفه
بايد وښيو چې که $z \in A$ وي، نو
$z \in A \cap (A \cup B)$۔] اوس [(1)] $z \in A$ فرض کړئ۔
بيا [(2)] $z \in A \cup B$ هم دے [ځکه له (1) څخه
$z \in A$ يا $z \in B$، چې د $\cup$ له تعريفه
د $z
  \in A \cup B$ معنا لري]؛ او ځکه $z \in A \cap (A \cup B)$
هم دے [ځکه د $\cap$ تعريف دواړه غواړي: $z \in A$، يعنې
(1)، او $z
  \in A \cup B)$، يعنې (2)۔ \emph{د سرچينې يادونه (OLSTH-041):
په دې وروستي چاپي عبارت کښې يو زيات تړونکے قوس شته؛
اصلي فورمول هماغسې ساتل شوے دے۔}]
\end{proof}

\begin{prob}
د $A \cup (A \cap B) = A$ لاندې ثبوت پراخ کړئ: ټولې
کارول شوې استنتاجي بڼې يادې کړئ، وښيئ چې هر ګام ولې
له فرضونو يا له مخکې ثابتو شوو دعوو راځي، او په کوم
ځاے کښې کوم تعريف ته اړتيا لرو۔
\begin{proof}
  که $z \in A \cup (A \cap B)$ وي، نو $z \in A$ يا
  $z \in A \cap B$۔ که $z \in A \cap B$ وي، $z \in A$۔
  هر $z \in A$، $\in A \cup (A
  \cap B)$ هم دے۔
\end{proof}
\end{prob}

\end{document}
